% !TEX root = ../elteikthesis_en.tex
\chapter{Introduction}
\label{ch:intro}

% ─────────────────────────────────────────────────────────────
\section{Motivation}
\label{sec:intro-motivation}

Students and staff at the Faculty of Informatics of Eötvös Loránd University face a recurring
practical challenge: finding accurate, up-to-date answers to questions about the curriculum.
Which courses must be completed before enrolling in a given subject? What are the rules for
repeating a failed exam? How does the Neptun system handle course registration? The answers to
these questions exist — scattered across faculty web pages, downloadable PDF regulations, DOCX
course descriptions, and administrative guidelines — but locating the right document, navigating
to the right section, and extracting a clear answer can take considerably longer than it should.

General-purpose search engines return ranked lists of pages; they do not synthesise an answer
from multiple sources or understand the specific terminology of Hungarian higher education. As
the volume of documentation grows each semester and more information migrates to sub-pages that
are hard to discover by browsing, the gap between the information that exists and the information
that students can efficiently find continues to widen.

% ─────────────────────────────────────────────────────────────
\section{Related Work}
\label{sec:intro-related}

Several universities have built AI-powered platforms that demonstrate the viability of language
models in higher education. Examining them critically, however, reveals a consistent set of
design assumptions that make them fundamentally different from the system developed in this
thesis.

\textbf{UM-GPT} (University of Michigan) is a university-wide AI platform running on Microsoft
Azure, offering access to multiple large language models to all students, faculty, and
staff~\cite{umich_genai}. Its companion tool, Maizey, is built on RAG and allows any user to
upload a custom dataset and create a personalised chatbot from it. While Maizey shares the same
underlying retrieval mechanism as the system in this thesis, it is a general-purpose tool: users
define their own knowledge base and their own purpose. The system developed here is
purpose-built — its knowledge base is fixed to official faculty documents and its goal is to
answer a specific and recurring student need, not to provide a blank canvas for arbitrary tasks.

\textbf{ZotGPT} (UC Irvine) is a broad AI services platform whose most relevant tool,
ClassChat, allows instructors to configure an AI assistant around their course syllabus and
materials~\cite{uci_zotgpt}. Its primary orientation is pedagogical — supporting teaching and
learning within a classroom context. It does not address the broader institutional information
problem: the rules, procedures, and regulations that exist outside any single course and that
students need to navigate independently.

\textbf{CreateAI} (Arizona State University) aggregates over 50 AI models into a single
platform for building custom AI applications across the entire
university~\cite{asu_createai}. Its defining characteristic is breadth, which is also its
limitation for the problem at hand: a platform built to support dozens of unrelated domains
simultaneously is not optimised for deep, precise retrieval over a single faculty's regulatory
documents.

Across all three systems, a common profile emerges: cloud infrastructure, general-purpose scope,
and significant IT investment. None are designed to answer a specific question set drawn from a
single document corpus, and none were evaluated on whether they actually answer domain-specific
questions correctly. This is the design space the present work inhabits.

% ─────────────────────────────────────────────────────────────
\section{Problem Statement}
\label{sec:intro-problem}

The related systems above establish that LLM-powered question answering is viable in higher
education. They do not, however, provide a usable template for the ELTE Faculty of Informatics.
Two concrete problems block a direct adoption.

\paragraph{Domain knowledge and hallucination.}
General-purpose LLMs have no reliable knowledge of ELTE-specific regulations, course codes, or
administrative procedures. Without grounding in official faculty documents, a language model
will hallucinate plausible-sounding but incorrect answers — precisely the failure mode that is
most harmful in an academic advising context, where a wrong answer about prerequisites or exam
rules can have real consequences for a student's academic progress.

\paragraph{Resource and infrastructure constraints.}
Deploying and maintaining a cloud-hosted platform at the scale of UM-GPT or CreateAI requires
dedicated IT teams and substantial recurring costs. A single faculty cannot replicate that
infrastructure. Any practical solution must run on commodity hardware — a standard laptop or a
modest server — with no GPU and no cloud subscription.

No dedicated, locally deployed question-answering system currently exists for the ELTE Faculty
of Informatics. Students and staff rely on direct web search, email queries to administrative
offices, or informal peer advice — all of which are slow, inconsistently accurate, and do not
scale as the faculty grows.

% ─────────────────────────────────────────────────────────────
\section{Objective}
\label{sec:intro-overview}

This thesis designs, implements, and evaluates a locally deployed question-answering system for
the ELTE Faculty of Informatics. The central objective is to demonstrate that a system built on
retrieval-augmented generation can provide students and staff with accurate, document-grounded
answers to questions about the curriculum and administrative procedures — without requiring the
infrastructure investment or general-purpose scope of the platforms examined in the related
work. The specific goals of the system are:

\begin{itemize}
  \item To collect and index publicly available faculty documents — including web pages, PDF
    regulations, and course descriptions — into a searchable local knowledge base that can be
    kept up to date as the faculty's web presence changes.

  \item To answer natural language questions about the curriculum, course requirements, and
    administrative procedures by retrieving relevant document passages and generating grounded
    responses, rather than relying on the general knowledge of a language model.

  \item To run entirely on commodity hardware without a GPU or cloud subscription, demonstrating
    that an accurate and responsive question-answering system does not require significant
    infrastructure investment.

  \item To provide a browser-based chat interface with source references, so that users can
    verify the origin of each answer and follow up on the source documents directly.

  \item To evaluate the system empirically against a no-RAG baseline, measuring whether
    retrieval meaningfully improves answer accuracy on questions drawn from real student queries.

  \item To build the system on a web service architecture that is deployment-ready, so that
    moving from a local prototype to a live university service would require infrastructure
    provisioning only, with no changes to the application code itself.
\end{itemize}

% ─────────────────────────────────────────────────────────────
\section{Technologies Used}
\label{sec:intro-technologies}

Two hard constraints shaped every technology decision: the system must run on commodity
hardware with 8~GB of RAM and no discrete GPU, and no query or document may leave the
local machine. The following components form the complete stack.

\begin{itemize}
    \item \textbf{Ollama}, a local model server, packages and serves large language
        models over a stable REST interface, handling quantisation and memory mapping
        automatically. It abstracts the complexity of model loading so that the rest
        of the application calls a single endpoint regardless of which model is active.
        Switching models requires only a one-line change in the configuration file,
        with no modifications to application code.

    \item \textbf{Llama~3.2~(3B) and Gemma~3~(4B)}, the two language models evaluated
        in this thesis, are open-weight models small enough to run on consumer hardware.
        Both are served at 4-bit quantisation, which reduces their memory footprint to
        approximately 2--2.5~GB and makes CPU-only inference feasible. This keeps the
        total memory usage well within the 8~GB RAM constraint while leaving headroom
        for the rest of the stack.

    \item \textbf{all-MiniLM-L6-v2}, a compact sentence transformer, encodes both
        document chunks and user queries into 384-dimensional dense vectors for semantic
        similarity search. Because the same model encodes both the index and the query,
        retrieval surfaces passages that are semantically close to the question rather
        than merely sharing keywords. It runs efficiently on CPU and provides a strong
        accuracy-to-resource trade-off for English factual retrieval tasks.

    \item \textbf{ChromaDB}, a lightweight vector database, stores the chunk embeddings
        and answers nearest-neighbour queries at runtime using cosine similarity. Unlike
        most vector databases, it runs in-process with no separate server and persists
        to a plain directory on disk. This satisfies the offline-operation requirement
        with no additional infrastructure to install or maintain.

    \item \textbf{FastAPI}, an asynchronous Python web framework, exposes the RAG
        pipeline as a REST API and serves the browser interface as static files from
        the same process. It was chosen for its low boilerplate, built-in request
        validation, and automatic API documentation. These properties make it practical
        for a single developer wrapping one core pipeline into a deployable HTTP service.

    \item \textbf{SQLite}, a serverless embedded database engine, logs every interaction
        — query, answer, sources, and response time — to a single file on disk with no
        server process required. It integrates directly into the Python application and
        requires no configuration beyond a file path. For a single-machine deployment
        with modest log volume, this is a simpler and more portable fit than a
        full database engine.

    \item \textbf{Vanilla JavaScript} powers the browser interface as a single
        self-contained HTML file with embedded CSS, served directly by the FastAPI
        backend. Avoiding a frontend framework eliminates a build step and all
        Node.js dependencies entirely. This keeps deployment as simple as placing
        one file in the static directory, consistent with the goal of minimal
        operational overhead.
        
    \item \textbf{Docker}, a containerisation platform, packages the entire
        application — the FastAPI backend, Ollama, and all dependencies — into
        isolated containers that run identically regardless of the host environment.
        This makes it the mechanism behind the deployability
        requirement: moving from a local machine to a faculty server requires
        only infrastructure provisioning, with no changes to the application code.
\end{itemize}

% ─────────────────────────────────────────────────────────────
\section{Thesis Structure}
\label{sec:intro-structure}

The remainder of this thesis is organised as follows.

\textbf{Chapter~\ref{ch:requirements}} defines the system requirements. It states
the functional requirements covering the capabilities exposed to the user, and the
non-functional constraints including resource limits, latency, privacy, and offline
operation.

\textbf{Chapter~\ref{ch:user}} provides user documentation: a use-case diagram, step-by-step
instructions for the web interface and command-line client, and annotated example queries.

\textbf{Chapter~\ref{ch:devdocs}} covers developer documentation. It describes the system
architecture and component design, then details the full implementation across the data
pipeline, embedding, RAG backend, REST API, logging, frontend, and deployment.

\textbf{Chapter~\ref{ch:testing}} presents the testing and evaluation of the system. It covers
functional testing, the empirical RAG vs.\ baseline evaluation, and a user evaluation.

\textbf{The Conclusion} summarises the findings, reflects on the key design decisions, and
outlines directions for future work.